% Supplementary computations, source pinpoints and certificate details.
% Bibliography keys refer to paper/paper.tex; numbered results and appendices below refer to the accompanying paper.

\section{The census for \texorpdfstring{$n\le8$}{n <= 8}}\label{sec:census}
Where an unordered collection of invariant factors is meant, it is called a multiset.

We computed the Smith form of $X_k$ on $S^\lambda_\Z$ for every
$\lambda\vdash n\le8$ and every $1\le k\le n$, 416 cases in all, and compared
it with Conjecture~7.4. The numbers of disagreements are:
\[
\begin{array}{l|cccccccc|c}
n & 1&2&3&4&5&6&7&8&\text{total}\\\hline
\text{cases} & 1&4&9&20&35&66&105&176&416\\
\text{disagreements} & 0&0&0&0&2&5&23&64&94
\end{array}
\]
The two disagreements at $n=5$ are those of Proposition~7.
Stanley's example $\lambda=(5,1)$, $k=5$, is reproduced: the Smith form is
$(1,1,3,3,12)$, as the survey states. Of the 94 disagreements, 23 have $k=3$
and are accounted for by Corollary~4. Disagreements involve the
prime $2$ in 79 cases and the prime $3$ in 30; they occur both for $k<n$ (85
cases) and for $k=n$ (9 cases), and for two-row shapes from $n=7$ on, the
first being $\lambda=(5,2)$.

The weaker remark fails first at $n=7$, for $(5,2)$ and $(2,2,1,1,1)$ with
$k=7$ (Proposition~2). Nor is the Smith form of $X_k$ always that
of the diagonal matrix of its eigenvalues: this holds in 279 of the 416
cases.

What survives in the data is limited: the conjecture holds for all
$n\le4$ (Proposition~7) and at $k=3$ for the shapes of
Corollary~4, and the Smith forms of $\lambda$ and $\lambda'$ agree
in every case. The modifications of the formula we tested all fail
(\emph{Other readings and modifications of the formula}), and we propose no corrected form.

\section{An open point}\label{sec:open}

The documented facts are these. Stanley reports computations through
$n\le7$, which he credits to Grinberg, and states that the conjecture was
made on the basis of these data; the survey defines $W_{\lambda,k}$ through a
basis of $\Z\sn$. Grinberg's later comment names the integral Specht module with its standard basis as the
intended lattice. On that lattice the conjecture fails at $n=5$
(Theorem~1), the weaker remark fails at $n=7$
(Proposition~2), and the census finds 30 disagreements with
$n\le7$. The data of
the computations reported in the survey are not available to us, and why
they supported the conjecture is not known to us. We attribute no error to these unavailable data.

\section{Verification}\label{sec:verify}

Three kinds of evidence are combined, and each result says which it has.
Lean proves statements about explicit algebraic objects: the identity of
Lemma~5; for the explicit matrices of
Theorem~1 and Proposition~2, their Smith forms, the
unreachability of the predicted forms, $AB=BA=2I$ for the first and
$A(A-4I)=0$ and the eigenvalues for the second; the matrix step of
Theorem~3; and the counting step of Corollary~4.
Standalone certificate checkers establish, by recomputation from the
definitions, that these matrices are those of $X_k$ on the integral Specht
lattices in the standard polytabloid basis, and that the predicted lists are
Stanley's; together with Lean this verifies Theorem~1 and
Proposition~2. A third checker verifies
Proposition~7 by exact computation over the 69 cases with
$n\le5$. Theorem~3, Corollary~4,
Proposition~8, Remark~6, the
exterior-square derivation in Appendix~D and the general statements about other
lattices are proved in writing and are not machine-checked; the census is
supporting data for them. The checks run with \file{checks/check.py} and
\texttt{lake build} in \file{lean/}. \emph{Scope of the checks} describes the
checkers, the census construction and the Lean statements, with what each
does not cover.


\section{Formal statements and computing details}\label{sec:formal-statements}
Lean verifies the following matrix and arithmetic statements. All declaration names below lie in the namespace \file{StanleySmithForms}. Here $M_d(\Z)$ denotes integer $d\times d$ matrices
and $\operatorname{diag}$ forms a diagonal matrix.
The integer matrix $A$ is displayed in Section~2; its fixed
unimodular matrices $U,V$ are displayed in Appendix~A. The theorem \file{counterexample} states exactly
\[
\begin{gathered}
(\det U,\det V\in\Z^\times\ \land\ UAV=\operatorname{diag}(1,1,1,1,2,2))\ \land\\
\neg\exists P,Q\in M_6(\Z):\quad
\det P,\det Q\in\Z^\times\ \land\ PAQ=\operatorname{diag}(1,1,1,1,1,4).
\end{gathered}
\]
Here $\Z^\times=\{1,-1\}$, so both determinants are units in $\Z$.
Two independent programs identify the matrix in Theorem~1:
$A$ represents $X_3$ in the ordered standard polytabloid basis, with column
$j$ the image of basis vector $j$, and the diagonal ruled out is the list conjectured in Theorem~1.
Lean does not construct that Specht module.

For the literal $14\times14$ matrices $A_{52},U_{52},V_{52}$ displayed here, the theorem \file{counterexample_52} states: $U_{52},V_{52}$
have unit determinants and $U_{52}A_{52}V_{52}$ is the diagonal listed in
Proposition~2; for every $\mu\in\Q$,
$(\exists v\in\Q^{14}\setminus\{0\}: A_{52}v=\mu v)$ if and only if
$\mu=0$ or $4$; and for every
$P,Q\in M_{14}(\Z)$ with unit determinants and every vector $d$,
$P\,A_{52}\,Q=\operatorname{diag}(d)$ implies some $d_i$ is divisible by $2$
and not by $4$. The declarations\allowbreak{} \file{eigen_identity_52} and
\file{charpoly_52} state, respectively, $A_{52}(A_{52}-4I)=0$ and
$\operatorname{charpoly}(A_{52})=t^5(t-4)^9$. Here $t$ is the polynomial indeterminate; the latter supplies the
multiplicities. The same column convention identifies $A_{52}$ with the matrix of $X_7$ in the standard polytabloid basis, and the programs check the conjectured list.

For every natural $n\ge3$, \file{group_ring_identity} states both
$X(X-T)=2$ and $(X-T)X=2$ in $\Z[\mathfrak S_n]$, where the transpositions defining $X$ and $T$ are $(0,2)+(1,2)$ and $(0,1)$ on $\{0,\dots,n-1\}$. These are the paper's letters
$1,2,3$ shifted by one; $2$ means $2$ times the ring identity.
The theorem \file{invariant_factors_of_mul_eq_two} states: for any finite
index set (with decidable equality in Lean), integer square matrices $A,B,U,V$ and
integer vector $d$, if $AB=2I$, $\det U,\det V$ are units and
$UAV=\operatorname{diag}(d)$, then each $d_i\in\{1,-1,2,-2\}$ and
$\#\{i:2\mid d_i\}=v_2(\det A)$, where $v_2$ is the $2$-adic valuation
of $|\det A|$.
It assumes a diagonalization; its entries may be signed and unordered.
The paper normalizes signs and divisibility order.

For natural $f,a,b$ with $a+b\le f$, \file{k3_counting} states
\[
\bigl(\forall r\in\{1,\dots,f\},\quad
2^{[r\le a]+[r\le b]}=
\begin{cases}2&r\le a+b,\\1&r>a+b\end{cases}\bigr)
\ \Longleftrightarrow\ (a=0\ \lor\ b=0),
\]
where brackets are $0$--$1$ indicators. Here $a,b$ count
tableaux with $3$ in $(1,3),(3,1)$, $f=f^\lambda$, and $r$ counts from the
largest factor, corresponding to the paper's index $f-r+1$.
In Theorem~3 and Corollary~4, the lattice passage from the group ring, the eigenvalue theorem, the count interpretations and the tableau extension/shape argument remain written. The census of \emph{The census for $n\le8$} agrees with both
results for $n\le8$: in all 63 cases with $k=3$ every invariant factor is $1$
or $2$, there are exactly $m$ factors equal to $2$, and the disagreements are
exactly the shapes with $\lambda_1\ge3$ and at least three rows.
A third Python checker establishes Proposition~7;
the census, an enumeration of all pairs $(\lambda,k)$ with $\lambda\vdash n\le8$ and $1\le k\le n$, is supporting data.


Lean~\cite{lean} and mathlib~\cite{mathlib} provide the formal proof
infrastructure. The formal proofs use propositional extensionality, quotient soundness and classical choice; the counting identity does not use classical choice. The computation also depends on the Python interpreter and
the two independent implementations of the Specht-matrix checkers.
The computation also reproduces Stanley's example $\lambda=(5,1)$, $k=5$, with Smith form $(1,1,3,3,12)$~\cite[example following Conjecture~7.4]{stanley}. Certificate checking took about two seconds on one CPU core. Reproduction instructions and the
pinned software versions are given in the accompanying README and
checking notes.

The recorded certificate check used one CPU core; its hardware model was not retained. The approximately two-second total is the retained certificate-check timing, not a new timing measurement. Full reproduction commands and software versions are in the README and the checker and Lean documentation.

\section{The conjecture in the paper's notation}
The matrix $W_{\lambda,k}$ is that
of $X_k$ in the standard polytabloid basis, and its Smith form is that of
$X_k$ as an endomorphism of $S^\lambda_\Z$, which does not depend on the
basis.
A rational eigenvalue description does not determine the integral Smith
form. Stanley's Conjecture~7.4~\cite{stanley} predicts
\[
\alpha_{f^\lambda-r+1}=\pm\prod_{(i,j)\in S_r}(j-i)
\qquad(1\le r\le f^\lambda),
\]
where $S_r$ consists of cells occupied by $k$ in at least $r$ standard
tableaux, and $j-i$ is the content of row $i$, column $j$.

\section{Sources and integral bases}\label{sec:sources}
The paper uses Stanley's Conjecture~7.4 and the following
paragraph~\cite{stanley}. The base lattice is the integral Specht module
specified in Grinberg's comment on Question~1~\cite{mo}.

An exact integral standard-basis reference is Grinberg's
Theorem~5.9.1~\cite{grinberg} (June 12, 2026 version, page~361), with base
ring $\Z$; his base-ring convention is on page~15. The content-spectrum
assertion used in Sections~3--4 and Appendix~C is in the paragraph preceding Conjecture~7.4 in Stanley's arXiv version~2, page~15~\cite{stanley}. For the standard-basis references above, the exact statements are
James~\cite[Theorem~8.4, page~29, and Corollary~8.6, page~30]{james}
and Sagan~\cite[Theorem~2.6.5(1), page~74]{sagan}, in the editions listed;
James's corollary gives integral spanning, whereas Sagan's displayed theorem
is stated over $\mathbb C$. For the rational content eigenbasis used below,
see Vershik and Okounkov~\cite[arXiv version~1, Corollary~2.10 and the following paragraph, page~13,
Theorem~5.8, page~21, and Proposition~6.1, page~22]{vo}; the second sentence after
that proposition's proof identifies the basis as Young's seminormal form.


The original Jucys and Murphy full texts were unavailable. For Murphy's page range we follow Grinberg's June 12, 2026 bibliography entry [Murphy81], which gives 287--297. Vershik and Okounkov's bibliography entry [24] instead gives 287--291; the original-text access gap leaves that discrepancy unresolved. The mathematical eigenvalue formulation used here is pinned to Vershik and Okounkov above; their Introduction, page~4, credits the independent rediscovery of the elements to Jucys and Murphy.

\section{Other readings and modifications of the formula}\label{app:readings}\label{sec:readings}

\paragraph{Other lattices for the counterexample (written; not machine-checked).} By
Theorem~3, the same Smith form arises on any lattice stable under
the permutations of $\{1,2,3\}$ in a module isomorphic to $S^{(3,1,1)}_\Q$:
for instance the lattices given by polytabloids built from column tabloids,
by a left ideal generated by a Young symmetrizer, by the dual lattice, or by
the saturated lattice $S^\lambda_\Q\cap M^\lambda_\Z$. No choice of lattice
or basis of this kind restores the prediction.


\paragraph{Other lattices in the census.} For $k=3$, Theorem~3 (proved in
writing) covers every lattice stable under the permutations of
$\{1,2,3\}$. Over the whole census the following were also compared. The
dual lattice
$\operatorname{Hom}(S^\lambda_\Z,\Z)$, with the contragredient action,
carries the transposed matrix, since transpositions are their own inverses;
its Smith form is the same. The saturated lattice
$S^\lambda_\Q\cap M^\lambda_\Z$, which contains $S^\lambda_\Z$, gives the
same Smith form as $S^\lambda_\Z$ in all 135 cases with $n\le6$; it was not
computed for larger $n$. The census also finds that $\lambda$ and its
conjugate $\lambda'$ have the same Smith form for every $k$, in all 416 cases,
as one expects from the classical relation between $S^{\lambda'}$ and the
dual of $S^\lambda$ twisted by the sign character~\cite[\S8]{james}, under
which $X_k$ becomes $-X_k$. The statements about the dual lattice and this
relation are written background, not machine-checked; the comparisons are
finite computations.

\paragraph{The literal group ring reading.} Read literally, the survey
compares the Smith form of left multiplication by $X_k$ on $\Z\sn$ with that
of a block diagonal matrix holding $f^\lambda$ copies of each predicted list.
Grinberg's comment explains why this is not the intended statement, and it
fails already for $n=4$, $k=4$. Writing $d^m$ for $m$ invariant factors equal
to $d$, left multiplication by $X_4$ on $\Z\mathfrak{S}_4$ has invariant factors
forming the multiset $\{0^4,1^{16},2^4\}$. The predicted lists are $(3)$ for $(4)$ and $(1^4)$,
$(1,2,2)$ for $(3,1)$ and $(2,1,1)$, and $(0,0)$ for $(2,2)$; with
multiplicities $f^\lambda$ the diagonal holds the multiset
$\{0^4,1^6,2^{12},3^2\}$, whose Smith normalization has invariant factors
forming the multiset $\{0^4,1^8,2^{10},6^2\}$. For $n\le3$ the two agree.

\paragraph{Modifications of the formula.}\label{sec:corrected}

We tested modifications of the formula against the census, the
Smith forms of $X_k$ on $S^\lambda_\Z$ for all $\lambda\vdash n\le8$. Write
$N_c$ for the number of standard tableaux of shape $\lambda$ in which $k$
occupies the cell $c$, so that the formula as stated uses
$S_r=\{c:N_c\ge r\}$. All sets of cells below consist of cells that $k$
occupies, that is, cells with $N_c\ge1$. The readings are:
\begin{itemize}\itemsep0pt
\item \emph{as stated}: $\alpha_{f-r+1}=|\prod_{c\in S_r}\mathrm{content}(c)|$;
\item \emph{multiset}: the Smith form of the diagonal matrix of these $f$
products, ignoring their order;
\item \emph{$\alpha_r$ indexing}: the $r$-th product assigned to $\alpha_r$
instead of $\alpha_{f-r+1}$;
\item \emph{exactly $r$, more than $r$, at most $r$}: $S_r$ replaced by
$\{c:N_c=r\}$, $\{c:N_c>r\}$ or $\{c:1\le N_c\le r\}$;
\item \emph{multiplicity classes}: one product for each positive value of
$N_c$, over the cells with that value, the list padded with $1$s to length
$f$ and assigned in increasing or decreasing order of $N_c$, or
Smith-normalized;
\item \emph{eigenvalue diagonal}: the Smith form of the diagonal matrix of
all eigenvalues of $X_k$ on $S^\lambda_\Q$;
\item \emph{zero contents dropped}: the formula as stated, with cells of
content $0$ left out of the products, except that a nonempty set consisting
only of cells of content $0$ gives $0$ (the empty set gives $1$);
\item \emph{$X_{k\pm1}$}: the formula for $k$ compared with the Smith form of
$X_{k+1}$ or of $X_{k-1}$;
\item \emph{$Y_k$}: the formula for $k$ compared with the Smith form of
$Y_k=\sum_{i>k}(i,k)$;
\item \emph{$\lambda'$}: the formula for the conjugate shape $\lambda'$
compared with $X_k$ on $S^\lambda$;
\item \emph{fillings of $1,\dots,k$}: $N_c$ replaced by the number of
distinct restrictions to the entries $1,\dots,k$ of standard tableaux of
shape $\lambda$ with $k$ in $c$;
\item \emph{weak form}: every nonzero invariant factor is a product of some
of the absolute values of the eigenvalues, the remark following the
conjecture.
\end{itemize}
A case matches when the predicted list equals the Smith form entry by entry.
The actual Smith form is written nonnegative, its nonzero entries in
increasing order and its zeros last. For the formula as stated, the
$\lambda'$, $X_{k\pm1}$ and $Y_k$ comparisons and the fillings reading, the
predicted list is put in the same order before the comparison. The ordered
trial lists of the other readings ($\alpha_r$ indexing, exactly, more than
or at most $r$, multiplicity classes, zero contents dropped) keep their
assigned positions $\alpha_1,\dots,\alpha_f$, zeros included; their
Smith-normalized versions are the Smith forms of the corresponding diagonal
matrices. Where a reading has an ordered and a Smith-normalized version, the
table gives the better count. The columns give
the matches among the 240 cases with $n\le7$ and the 416 with $n\le8$ (196
and 350 for $X_{k\pm1}$, where $k\pm1$ must lie in $1,\dots,n$), whether the
reading reproduces Stanley's example $(1,1,3,3,12)$ for $\lambda=(5,1)$,
$k=5$, and the first case, in order of $n$, where it fails.
\[
\begin{array}{lccll}
\text{reading} & n\le7 & n\le8 & (5,1),\,k=5 & \text{first failure}\\\hline
\text{as stated} & 210 & 322 & \text{yes} & (3,1,1),\ k=3\\
\text{multiset} & 210 & 322 & \text{yes} & (3,1,1),\ k=3\\
\lambda' & 210 & 322 & \text{yes} & (3,1,1),\ k=3\\
\alpha_r\text{ indexing} & 122 & 181 & \text{no} & (3,1),\ k=3\\
\text{exactly }r & 96 & 133 & \text{no} & (2,1),\ k=1\\
\text{more than }r & 47 & 70 & \text{no} & (1),\ k=1\\
\text{at most }r & 94 & 131 & \text{no} & (2,1),\ k=1\\
\text{multiplicity classes} & 96 & 133 & \text{no} & (2,1),\ k=1\\
\text{eigenvalue diagonal} & 177 & 279 & \text{yes} & (3,2),\ k=4\\
\text{zero contents dropped} & 173 & 266 & \text{yes} & (3,2),\ k=4\\
X_{k+1} & 3 & 3 & \text{no} & (2),\ k=1\\
X_{k-1} & 3 & 3 & \text{no} & (2),\ k=2\\
Y_k & 26 & 26 & \text{no} & (2),\ k=1\\
\text{fillings of }1,\dots,k & 144 & 203 & \text{yes} & (2,1),\ k=1\\
\text{weak form} & 238 & 410 & \text{yes} & (5,2),\ k=7
\end{array}
\]
The multiset reading agrees with the formula as stated in every case, since
$S_1\supseteq S_2\supseteq\cdots$ makes the stated products a divisibility
chain. The $\lambda'$ reading agrees with it as well: conjugation moves the
counts to transposed cells and changes the sign of the contents, and the
census finds the same Smith form for $\lambda$ and $\lambda'$. So does the
reading with the labels counted from $n$ downward, $Y_{n+1-k}=w_0X_kw_0$ with
$w_0$ the reversal permutation, whose Smith form equals that of $X_k$ in all
416 cases.

The modifications we tested all fail on the data with $n\le7$. The best fit is the formula
as stated, together with its equivalent readings; it matches 210 of the 240
cases with $n\le7$ and first fails at $\lambda=(3,1,1)$, $k=3$. Every other
reading that predicts the whole Smith form first fails at $n\le5$, most of
them at $n\le3$. The weak form, which predicts only a property of it, holds
for all $n\le6$ and fails at $n=7$, for $(5,2)$ and $(2,2,1,1,1)$ with $k=7$
(Proposition~2). We propose no corrected form. The program for
this comparison is included with the optional census programs; it needs the
python-flint library.


\section{Scope of the checks}\label{app:checks}

\paragraph{Certificates and checkers.} The two counterexamples,
$\lambda=(3,1,1)$ with $k=3$ and $\lambda=(5,2)$ with $k=7$, are given as
certificates, each holding the standard tableaux in order, the
matrix $A$ of $X_k$, unimodular $U$ and $V$ with $UAV$ diagonal, the Smith
form and the prediction. Two standalone checkers, written independently and
using only the Python standard library, read them; neither trusts the
encoded matrix. The first enumerates the standard
tableaux from all $n!$ fillings, proves that the standard polytabloids form a
$\Z$-basis of the $\Z$-span of all $n!$ polytabloids, recomputes $A$ exactly
from the tabloid module, checks $UAV=\operatorname{diag}(\text{snf})$ with
integer $U,V$ of determinant $\pm1$ and a divisibility chain, and recomputes
the prediction. The second, with separate code, enumerates the standard
tableaux by lattice words, shows that the $\F_2$-span of all polytabloids has
dimension $f^\lambda$, computes the rank of $X_k$ there without reference to
a basis, compares it with the number of odd predicted factors, verifies
$X_ke_{T_j}=\sum_iA_{ij}e_{T_i}$ in the tabloid module for every $j$, and
checks $U$, $V$ and the Smith form with its own arithmetic. For a
certificate that also claims the failure of the weak remark, both checkers,
each with its own code, take as eigenvalues the contents of the cell holding
$k$ in each standard tableau; verify $\prod_c(A-cI)=0$ over the distinct
contents $c$, so that $A$ is diagonalizable with eigenvalues among them;
check $\operatorname{tr}(A^j)=\sum c^j$ over the eigenvalue list for $1\le j$
below the number of distinct contents, which fixes the multiplicities; and
require some Smith entry that is not the absolute value of a product of a
sub-multiset of the eigenvalues (the empty product being $1$). The checkers are independent: separate code, different
enumerations of the standard tableaux (all $n!$ fillings, and lattice
words), and different arithmetic (exact rational elimination, and Bareiss
determinants with linear algebra over $\F_2$). Both reject eight corrupted certificates (a changed entry of $A$; the
Smith form replaced by the prediction; a non-unimodular $U$; a rational $U$
of determinant $\pm1$; an extra factor; a non-standard basis tableau; an
altered prediction; and, in the weak-form certificate, a planted change to
$A$ with $U$ adjusted so that $UAV$ is still the claimed diagonal matrix,
rejected because $A(A-4I)\ne0$). Both also reject three controls: the
agreeing case $\lambda=(5,1)$, $k=5$, as showing no discrepancy, and two
weak-form claims about true statements, $(5,1)$ with $k=5$, where
$12=4\cdot3$ is a product of eigenvalues, and $(3,1,1)$ with $k=3$, where
$2=2\cdot1$ is one. A third checker, written separately from both and
importing neither, recomputes every $\lambda\vdash n\le5$ and every $k$: it
proves that the standard polytabloids form a $\Z$-basis, forms the matrix of
$X_k$, computes its Smith form from determinantal divisors, and asserts
Proposition~7 and Stanley's example $\lambda=(5,1)$, $k=5$.
The entry point is \file{checks/check.py}, with the certificates in
\file{data/}.

\paragraph{Census.} The census was computed by two constructions. The first
generates $S^\lambda_\Z$ inside $M^\lambda_\Z$ as the $\Z\sn$-span of one
polytabloid, closing under simple transpositions with Hermite normal forms;
it uses neither the standard basis nor straightening, and every Smith form
is computed by two implementations. The second uses the standard polytabloid
basis and exact rational linear algebra. They agree in all 135 cases with
$n\le6$. The census output is included as \file{data/census-n8.json}; the
census programs are optional and not part of the check.

\paragraph{Lean.}\sloppy The checked statements are in Lean~4~\cite{lean} with
mathlib~\cite{mathlib}, in the namespace \texttt{StanleySmithForms}.
\begin{itemize}\itemsep0pt
\item \texttt{group\_ring\_identity}: the identity of
Lemma~5 in the integral group ring of $\sn$, for every
$n\ge3$.
\item \texttt{A\_det}: $\det A=-4$; \texttt{U\_det}, \texttt{V\_det}:
$\det U=-1$, $\det V=1$; \texttt{smith\_form}:
$UAV=\operatorname{diag}(1,1,1,1,2,2)$.
\item \texttt{A\_mul\_B}, \texttt{B\_mul\_A}: $AB=BA=2I$ for the explicit
integer matrix $B=2A^{-1}$.
\item \texttt{rank\_mod\_two}: $A$ has rank $4$ over $\F_2$;
\texttt{not\_predicted}: no integer matrices $P,Q$ with unit determinant
satisfy $PAQ=\operatorname{diag}(1,1,1,1,1,4)$; \texttt{counterexample}
combines the Smith form and the unreachability of the predicted form.
\item \texttt{invariant\_factors\_of\_mul\_eq\_two}: for integer matrices
$A,B$ with $AB=2I$, every unimodular diagonalization $UAV=\operatorname{diag}(d)$
has $d_i\in\{\pm1,\pm2\}$, and the number of even $d_i$ is $v_2(\det A)$
(applying it to the nonnegative Smith form of an arbitrary lattice is the
written step of Theorem~3); \texttt{witness\_factors} applies it
to the witness.
\item For $\lambda=(5,2)$, $k=7$, with the explicit $14\times14$ matrix $A$:
\texttt{smith\_form\_52}, \texttt{U52\_unimodular},
\texttt{V52\_unimodular}: the Smith form via unimodular $U$, $V$;
\texttt{eigen\_identity\_52}: $A(A-4I)=0$; \texttt{charpoly\_52}: the
characteristic polynomial of $A$ is $x^5(x-4)^9$; \texttt{eigenvalues\_52}:
the rational eigenvalues of $A$ are exactly $0$ and $4$;
\texttt{weak\_remark\_fails\_52}: every unimodular diagonalization of $A$
has an entry divisible by $2$ but not by $4$, proved from a parity certificate,
integer vectors $x,y,w,z$ with $Ay=2x$, $w^{\mathsf T}A=2z^{\mathsf T}$ and
$w\cdot x=1$; \texttt{counterexample\_52} combines the Smith form, the
eigenvalues and this statement. The identities between explicit
$14\times14$ matrices are checked by evaluation in the Lean kernel, which
adds no axioms.
\item \texttt{k3\_counting}: the counting step of Corollary~4, as
stated in Section~4.
\end{itemize}
Lean does not construct Specht modules, Young's seminormal basis or standard
tableaux. It does not check that the explicit matrices are those of $X_3$ on
$S^{(3,1,1)}_\Z$ and of $X_7$ on $S^{(5,2)}_\Z$, nor that the predicted lists
and the contents are Stanley's; the certificate checkers, which recompute
these from $\lambda$ and $k$, cover that step for both counterexamples. It does not check the passage from the group ring to an
arbitrary lattice, the content spectrum, or the facts about tableaux used in
Theorem~3 and Corollary~4; those are written.\par\fussy




\paragraph{The shape $(5,2)$ at $k=7$.}
Here $N_c$ denotes the number of standard tableaux of shape $(5,2)$ in
which $7$ occupies the cell $c$. We have $N_{(1,5)}=9$ and $N_{(2,2)}=5$.

\paragraph{The finite calculation through $n\le5$.}
The cases are all 69 pairs $(\lambda,k)$ with $\lambda\vdash n\le5$ and $1\le k\le n$. The standard basis property means that the standard polytabloids form a $\Z$-basis of their span. The program constructs these polytabloids from the definitions and establishes this property in every case. It computes Smith forms using determinantal divisors: $d_j$ is the gcd of all $j\times j$ minors, computed as exact integer determinants, and the $j$-th invariant factor is $d_j/d_{j-1}$. The separate enumeration through $n\le8$ agrees. The method is detailed in checks/README.md.
